\documentclass[11pt,a4paper]{article}

\usepackage[T1]{fontenc}
\usepackage[utf8]{inputenc}
\usepackage[ngerman]{babel}
\usepackage{lmodern}
\usepackage{microtype}
\usepackage{amsmath,amssymb,amsthm,mathtools}
\usepackage{mathrsfs}
\usepackage[a4paper,margin=2.7cm]{geometry}
\usepackage{graphicx}
\usepackage{xcolor}
\usepackage{hyperref}
\usepackage{enumitem}
\usepackage{booktabs}

\newtheorem{definition}{Definition}[section]
\newtheorem{satz}[definition]{Satz}
\theoremstyle{remark}
\newtheorem{bemerkung}[definition]{Bemerkung}
\newtheorem{hypothese}[definition]{Hypothese}

\newcommand{\N}{\mathbb{N}}
\newcommand{\R}{\mathbb{R}}
\newcommand{\C}{\mathbb{C}}
\newcommand{\T}{\mathbb{T}}

\title{Die Uhr in den Schleifen\\
	\large Primzahruhr, Holonomie und Feinstruktur im Bamberg-Modell}
\author{Thomas Hoffbauer\\Bamberg, Deutschland}
\date{\today}

\begin{document}
	\maketitle
	
	\begin{abstract}
		Dieser Text ist kein klassischer Fachaufsatz, sondern ein Positionsartikel. Er verfolgt die These, dass in einem diskreten Schalenmodell arithmetische Familien, geometrische Schleifen und spektrale Feinstrukturen auf eine gemeinsame Quelle zurückgeführt werden können: auf eine Uhr, die nicht auf Ziffernblättern, sondern auf Kanten, Plaquetten und Phasen läuft. Ausgangspunkt sind Tetraederschichten und eine gewichtete Primzahruhr. Daraus wird ein endlicher Gauge-Dirac-Operator auf einem diskreten Kanalraum aufgebaut, in dem echte Aharonov--Bohm-Flüsse auf Schleifen wirksam werden. Ein Feinstruktur-Shift der Form $2\pi/137$ erweist sich dabei nicht als beliebige Zugabe, sondern als wirksame Feinjus­tierung der Holonomie. In einer periodisierten Minimalversion treten Berry-Krümmung und nichttriviale Chern-Zahlen auf. Der Text versteht diese Resultate als Hinweis darauf, dass zwischen diskreter Geometrie, Primzahlordnung und physikalischer Phasenstruktur mehr Verwandtschaft besteht, als in üblichen Beschreibungen sichtbar wird.
	\end{abstract}
	
	\section{Worum es hier eigentlich geht}
	
	Es gibt mathematische Bilder, die zunächst wie freie Metaphern wirken, bis man entdeckt, dass sie numerisch zu arbeiten beginnen. Genau in diesem Zwischenraum bewegt sich der vorliegende Text. Sein Gegenstand ist eine Uhr, aber nicht im gewöhnlichen Sinn. Sie misst nicht Minuten, sondern Verschiebungen von Primfamilien, Spuren von Holonomien, Intensitäten von Kanalfiltern und die topologische Reaktion eines Operators auf Flüsse.
	
	Der Ausgangspunkt ist bescheiden: natürliche Zahlen werden in Tetraederschichten organisiert. Darauf erscheinen glatte Zahlen, Defektzahlen und Primzahlen nicht als formloser Strom, sondern als strukturierte Besetzung äußerer Hüllen. Zunächst könnte man das als bloße arithmetische Spielerei ansehen. Doch sobald man die Primfamilien modulo $12$ gewichtet, aus den drei aktiven Familien $A,B,C$ eine Uhr baut und diese Uhr nicht nur betrachtet, sondern in die Kanten eines diskreten Operators hineingibt, beginnt das Bild, numerisch ernst zu werden.
	
	Die zentrale Einsicht ist dabei einfach zu sagen und schwer zu ersetzen:
	\begin{quote}
		Die relevante Information sitzt im Modell nicht primär in diagonalen Werten, sondern auf Kanten und Schleifen.
	\end{quote}
	
	Darum ist es kein Zufall, dass ein bloßer diagonaler Uhrterm wenig hilft, ein Kantenfilter aber schon, und ein echter Plaquettenfluss noch viel mehr. Der Weg führt von der Zahl zur Kante, von der Kante zur Schleife, von der Schleife zur Holonomie.
	
	\section{Tetraederschichten: die erste Uhr}
	
	Jede Uhr braucht ein Zifferblatt. In diesem Modell übernehmen Tetraederschichten diese Rolle. Die Tetraederzahlen
	\begin{equation}
		T_L=\frac{L(L+1)(L+2)}{6}
	\end{equation}
	legen diskrete Hüllen fest. Zwischen zwei aufeinanderfolgenden Tetraederzahlen liegt eine Schale, und diese Schale wird von Zahlen besetzt. In dieser einfachen Konstruktion erscheint bereits der erste Rhythmus: die Schichtparität. Ungerade und gerade Schichten verhalten sich nicht gleich. Die einen lassen sich als Tri-Effekt lesen, die anderen als oktonischer Defekt. Ob man diese Begriffe wörtlich oder heuristisch nimmt, ist zunächst zweitrangig. Wichtig ist, dass die Parität numerisch nicht völlig belanglos bleibt.
	
	Auf diesen Hüllen lebt die Primzahruhr. Primzahlen der Restklassen $1,5,7,11$ modulo $12$ werden nicht bloß gezählt, sondern logarithmisch gewichtet. Gerade diese Gewichtung macht aus einer groben Uhr eine feine Uhr. Die Familien tragen nicht nur Häufigkeiten, sondern Lasten. Damit wird die Uhr nicht nur kombinatorisch, sondern energetisch lesbar.
	
	\section{Warum die Uhr nicht diagonal sein darf}
	
	Es wäre naheliegend, die so erhaltene Uhramplitude einfach als zusätzlichen Energieterm auf die Zustände zu legen. Genau das funktioniert jedoch schlecht. Der numerische Befund ist eindeutig: Als diagonale Energie ruiniert die Uhramplitude eher die spektrale Feinstruktur, als dass sie sie verbessert. Das ist ein wichtiges Negativresultat.
	
	Denn daraus folgt etwas Positives. Die Primzahruhr ist keine Potentialhöhe. Sie ist eine Auswahlregel. Sie reguliert, welche Übergänge bevorzugt werden und welche nicht. Sie gehört damit in die Leitfähigkeit des Modells, nicht in dessen statische Selbstbeschreibung.
	
	Dies ist vielleicht die erste wirkliche physikalische Pointe des Bildes. Ein System, dessen relevanteste Größe nicht in den Punkten, sondern in den Übergängen sitzt, nähert sich dem, was man von Gauge-Physik erwartet.
	
	\section{Der diskrete Dirac-Operator als eigentliches Organ}
	
	Sobald der E0/ABC-Split eingeführt ist, besitzt jeder Zustand einen isotropen Kern und eine anisotrope Excitation. Daraus lässt sich ein anisotroper diskreter Dirac-Operator aufbauen. Seine Koeffizienten folgen den drei Kanalrichtungen, seine Differenzen tragen komplexe Kantenphasen, und sein Massenterm bleibt vergleichsweise mild. Das Entscheidende ist: Der Operator lebt nicht von einer einzigen Größe, sondern von einem fein austarierten Zusammenspiel aus Kern, Anisotropie, Rough-Komponente und Uhrfilter.
	
	Gerade hier zeigt sich, dass die bessere Sprache des Modells nicht die des Potentialtopfes, sondern die des Transports ist. Zustände sind nicht primär Speicherplätze, sondern Knoten eines gerichteten, phasensensitiven Übertragungsnetzes.
	
	\section{Holonomie: wo das Modell ernst wird}
	
	Bis zu diesem Punkt könnte man noch behaupten, alles sei eine geschickte geometrische Umverpackung. Das ändert sich in dem Moment, in dem echte Schleifenflüsse eingebaut werden. Ein globaler Phasenoffset ist harmlos; er kann oft weggeeicht werden. Ein echter Aharonov--Bohm-Fluss durch elementare Plaquetten ist etwas anderes. Er hinterlässt eine nichttriviale Holonomie, und damit eine reale Differenz zwischen Wegen.
	
	Numerisch ist dieser Schritt der Wendepunkt. Mit echtem Plaquettenfluss verbessert sich die lokale Feinstruktur des Spektrums deutlich. Das Modell reagiert nicht auf die bloße Existenz eines Phasenparameters, sondern auf dessen Schleifenwirksamkeit. Dies ist von zentraler Bedeutung, weil es den geometrischen Grundgedanken bestätigt: Nicht die bloße Lage eines Zustands, sondern die Phasenakkumulation entlang geschlossener Wege ist der eigentliche Träger der Information.
	
	\section{Die Feinstrukturkonstante als mikroskopischer Holonomie-Shift}
	
	Ein besonders aufschlussreicher Befund betrifft die Zahl $1/137$. Als nackte Zusatzphase auf einem einzelnen Kanal bewirkt sie fast nichts. In dieser Form bleibt sie zu äußerlich. Anders verhält es sich, wenn sie den echten Plaquettenfluss verschiebt:
	\begin{equation}
		\Phi \mapsto \Phi + \frac{2\pi}{137}.
	\end{equation}
	
	Dann verbessert sich die numerische Feinstruktur im Modell drastisch. Die Spacing-Korrelation steigt in der untersuchten Konfiguration von ungefähr $0.322$ auf ungefähr $0.562$. Das ist zu deutlich, um als bloße kosmetische Modifikation abgetan zu werden.
	
	Was bedeutet das? Nicht, dass die Feinstrukturkonstante im Modell bereits physikalisch erklärt wäre. Aber es bedeutet, dass sie in dieser Architektur an der richtigen Stelle sitzt, wenn sie als mikroskopische Holonomie-Justierung verstanden wird. Nicht als rohe Zahl an irgendeinem Kanal, sondern als feiner Versatz der Schleifenphase.
	
	Dies ist eine der stärksten Einsichten des gesamten Projekts. Die Feinstruktur erscheint nicht als Eigenschaft eines isolierten Punktes, sondern als Qualität einer interferierenden Schleifenwelt.
	
	\section{Euler, Aharonov, Chern}
	
	Mit dem Eintritt der Schleifen stellt sich sofort die Frage nach der Topologie. Hier muss man sauber unterscheiden. Die Eulerzahl des diskreten Trägers beschreibt die kombinatorische Form des zugrunde liegenden Komplexes. Im untersuchten Fall ergibt sich für den dreidimensionalen Zellkomplex die Zahl
	\begin{equation}
		\chi=1.
	\end{equation}
	
	Das spricht für einen zusammenhängenden Kernkörper. Diese Zahl gehört zur Geometrie des Trägers.
	
	Etwas anderes ist die Aharonov--Bohm-Holonomie. Sie beschreibt die lokale Phasenakkumulation auf Schleifen. Und noch einmal etwas anderes ist die Chern-Zahl eines periodisierten Spektralbündels. Sie gehört nicht dem Raum selbst, sondern dem Band des Operators auf diesem Raum.
	
	Gerade die Klarheit dieser Trennung ist wichtig. Die Eulerzahl sagt, welche Art Raum wir haben. Die Holonomie sagt, wie Phasen sich auf Schleifen verhalten. Die Chern-Zahl sagt, in welcher globalen topologischen Phase das periodisierte Spektrum lebt.
	
	\section{Von Klitzing: das globale Echo}
	
	Sobald das Modell periodisiert wird, treten Berry-Krümmung und Chern-Zahl auf. Damit erscheint auch das vom-Klitzing-Motiv in seiner richtigen Rolle: nicht als rhetorische Verzierung, sondern als globale, quantisierte Antwort auf eine lokal aufgebaute Holonomiestruktur.
	
	Der interessante Befund ist, dass der numerisch beste feinstrukturverschobene Fluss nicht in einer vollständig trivialen topologischen Umgebung liegt. Er befindet sich vielmehr in der Nachbarschaft nichttrivialer Chern-Plateaus. Das heißt nicht, dass der beste Nullstellenfluss selbst schon maximal topologisch robust wäre. Aber es heißt, dass er in eine Zone fällt, in der topologische Phasenwechsel greifbar werden.
	
	Dieses Resultat ist vielleicht das überzeugendste Bindeglied zwischen dem endlichen Nullstellenmodell und dem periodischen Bandmodell. Beide reagieren auf denselben Flussparameter, aber auf unterschiedliche Weise: das eine lokal-spektroskopisch, das andere global-topologisch.
	
	\section{Die drei Geschwindigkeiten als offener Gedanke}
	
	Man kann an dieser Stelle noch einen weiteren Gedanken anknüpfen. Wenn die drei Kanäle $A,B,C$ nicht nur Richtungen, sondern effektive Transportmoden sind, dann liegt es nahe, ihnen drei verschiedene Geschwindigkeiten $v_A,v_B,v_C$ zuzuschreiben. Diese müssten nicht mit beobachtbaren Wechselwirkungskonstanten identisch sein, könnten aber an drei verschieden starke Kopplungsregime erinnern. Die beobachtbare Energiegeschwindigkeit wäre dann keine primitive Einzelgröße, sondern eine kohärente Vereinigungsgröße der drei Kanalgeschwindigkeiten.
	
	Dieser Gedanke ist in der vorliegenden Fassung noch heuristisch. Er fügt sich jedoch organisch in das Modell ein, da die periodisierte 2-Band-Version bereits mit anisotropen Geschwindigkeiten arbeitet. Die Lichtgeschwindigkeit würde in dieser Lesart nicht als Eigenschaft eines einzelnen Kanals, sondern als Grenzgeschwindigkeit kohärenter Vereinigung erscheinen.
	
	\section{Was dieses Modell nicht behauptet}
	
	Gerade weil die numerischen Befunde interessant sind, ist eine methodische Zurückhaltung wichtig. Das Modell behauptet nicht, die Riemann-Nullstellen erklärt zu haben. Es behauptet auch nicht, die Feinstrukturkonstante physikalisch hergeleitet zu haben. Und es behauptet schon gar nicht, das Standardmodell in ein paar diskrete Kanalphasen zu übersetzen.
	
	Was es behauptet, ist bescheidener und vielleicht gerade deshalb ernst zu nehmen: Es gibt eine numerisch nichttriviale Architektur, in der diskrete Schalen, Primfamilien, Schleifenholonomien, feine Flussverschiebungen und topologische Bandphasen miteinander kommunizieren. Diese Architektur ist reich genug, um echte Verbesserungseffekte zu tragen, und einfach genug, um noch lesbar zu bleiben.
	\begin{table}[t]
		\centering
		\caption{Numerische Befunde zum Eskalationsmodell auf Tetraederschichten. Die Werte fassen die wichtigsten Tests für Tri-/Okto-Regime sowie die Schwellen $5$, $7$ und $11$ zusammen.}
		\label{tab:eskalation}
		\begin{tabular}{@{}llr@{}}
			\toprule
			Bereich & Größe & Wert \\
			\midrule
			Tri-/Okto-Regime
			& $\overline{P_{\mathrm{Tri}}}\,$ auf Tri-Schichten & $0.04086$ \\
			Tri-/Okto-Regime
			& $\overline{P_{\mathrm{Tri}}}\,$ auf Okto-Schichten & $0.02312$ \\
			Tri-/Okto-Regime
			& $\overline{P_{\mathrm{Okto}}}\,$ auf Tri-Schichten & $0.86039$ \\
			Tri-/Okto-Regime
			& $\overline{P_{\mathrm{Okto}}}\,$ auf Okto-Schichten & $0.90481$ \\
			Tri-/Okto-Regime
			& mittlere Defektdichte auf Tri-Schichten & $0.74477$ \\
			Tri-/Okto-Regime
			& mittlere Defektdichte auf Okto-Schichten & $0.79745$ \\
			\midrule
			Umschläge
			& mittlerer Defektsprung $\,\Delta \rho_{\mathrm{defect}}$ & $0.04658$ \\
			Umschläge
			& Anteil positiver Defektsprünge & $0.61905$ \\
			Umschläge
			& Anteil von Familienwechseln & $0.80952$ \\
			\midrule
			Schwelle $5$
			& $\mathrm{corr}(P_5,\mathrm{family\_change}_{\mathrm{next}})$ & $-0.43330$ \\
			\midrule
			Schwelle $7$
			& $\mathrm{corr}(\mathrm{core7},\mathrm{family\_change}_{\mathrm{next}})$ & $0.30722$ \\
			Schwelle $7$
			& $\mathrm{corr}(\pi_7,\mathrm{family\_change}_{\mathrm{next}})$ & $-0.50637$ \\
			Schwelle $7$
			& $\mathrm{corr}(\rho_7,\Delta \rho_{\mathrm{defect,next}})$ & $-0.14455$ \\
			\midrule
			Schwelle $11$
			& $\mathrm{corr}(\mathrm{core11},\mathrm{family\_change}_{\mathrm{next}})$ & $0.47417$ \\
			Schwelle $11$
			& $\mathrm{corr}(\rho_{11},\Delta \rho_{\mathrm{defect,next}})$ & $-0.52278$ \\
			Schwelle $11$
			& $\mathrm{corr}(\rho_{11},\Delta |\zeta_{\mathrm{clock}}|_{\mathrm{next}})$ & $-0.45039$ \\
			\bottomrule
		\end{tabular}
	\end{table}
	
	\section{Schluss}
	
	Der eigentliche Gewinn dieses Modells liegt vielleicht nicht in einer einzigen Formel, sondern in einer Perspektive. Zahlen erscheinen nicht mehr als bloße Positionen auf einer Geraden, sondern als Besetzung diskreter Schalen. Primzahlen erscheinen nicht mehr nur als isolierte Punkte, sondern als gewichtete Familien auf einer Uhr. Spektrale Information sitzt nicht mehr in Potentialen, sondern in Kanten und Schleifen. Und die Feinstrukturkonstante wirkt nicht mehr als zufällige Naturzahl, sondern als feiner Versatz einer Holonomie.
	
	So entsteht ein Bild, in dem Euler, Aharonov und von Klitzing nicht nebeneinander stehen, sondern drei Schichten derselben Geschichte erzählen: die Gestalt des Trägers, die Phase auf der Schleife und die globale Klasse des Spektrums.
	
	\begin{thebibliography}{99}
		\bibitem{AharonovBohm}
		Y. Aharonov, D. Bohm,
		Significance of electromagnetic potentials in the quantum theory,
		\emph{Phys. Rev.} \textbf{115} (1959), 485--491.
		
		\bibitem{Klitzing}
		K. von Klitzing, G. Dorda, M. Pepper,
		New method for high-accuracy determination of the fine-structure constant based on quantized Hall resistance,
		\emph{Phys. Rev. Lett.} \textbf{45} (1980), 494--497.
		
		\bibitem{Thouless}
		D. J. Thouless, M. Kohmoto, M. P. Nightingale, M. den Nijs,
		Quantized Hall conductance in a two-dimensional periodic potential,
		\emph{Phys. Rev. Lett.} \textbf{49} (1982), 405--408.
		
		\bibitem{Fukui}
		T. Fukui, Y. Hatsugai, H. Suzuki,
		Chern numbers in discretized Brillouin zone: Efficient method of computing (spin) Hall conductances,
		\emph{J. Phys. Soc. Japan} \textbf{74} (2005), 1674--1677.
		
		\bibitem{Nakahara}
		M. Nakahara,
		\emph{Geometry, Topology and Physics},
		2nd ed., Taylor \& Francis, Boca Raton, 2003.
	\end{thebibliography}
	
\end{document}
